% part: intuitionistic-logic
% chapter: soundness-completeness
% section: soundness-nd

\documentclass[../../../include/open-logic-section]{subfiles}

\begin{document}

\olfileid{int}{sc}{snd}

\olsection{సహజ నిగమనం నిర్దుష్టత}

ఇప్పుడు సంబంధ అర్థవిజ్ఞానానికి సంబంధించి సహజ నిగమనం
నిర్దుష్టతను నిరూపిస్తాం. అంటే, పరికల్పనల సమితి నుంచి
\tetoken{ఒక సూత్రం}{!!a{formula}} \tetoken{వ్యుత్పాదించదగినదైతే}{!!{derivable}},
ఆ పరికల్పనల సమితి ఆ \tetoken{సూత్రాన్ని}{!!{formula}}
అనుగమింపజేస్తుందని చూపుతాం.

\begin{thm}[నిర్దుష్టత]
  \ollabel{thm:soundness}
  $\Gamma \Proves !A$ అయితే, $ \Gamma \Entails !A$.
\end{thm}

\begin{proof}
  $\Gamma \Proves !A$ అయితే $\Gamma \Entails !A$
  అని నిరూపిస్తాం. $\Gamma$ నుంచి~$!A$ యొక్క
  \tetoken{వ్యుత్పత్తి}{!!{derivation}}పై ఆగమన
  పద్ధతిలో నిరూపణ సాగుతుంది.

  \begin{enumerate}
  \item \tetoken{వ్యుత్పత్తి}{!!{derivation}} కేవలం
    పరికల్పన~$!A$తోనే ఉంటే, $!A \Proves !A$;
    $!A \Entails !A$ అని చూపాలి.
    $\mSat{M}{!A}[w]$ అనుకుందాం. అప్పుడు సహజంగానే
    $\mSat{M}{!A}[w]$.

  \item \tetoken{వ్యుత్పత్తి}{!!{derivation}} చివర
    $\Intro{\land}$ ఉంటే:
    \iftag{probAnd}{అభ్యాసం.}{\tetoken{ఉపసంహరించని}{!!{undischarged}}
      పరికల్పనలైన~$\Gamma$ నుంచి $!B$ పూర్వాపేక్షకు,
      \tetoken{ఉపసంహరించని}{!!{undischarged}}
      పరికల్పనలైన~$\Delta$ నుంచి~$!C$ పూర్వాపేక్షకు
      ఉన్న \tetoken{వ్యుత్పత్తులు}{!!{derivation}s}
      $\Gamma \Proves !B$, $\Delta \Proves !C$ అని
      చూపుతాయి. ఆగమన పరికల్పన ప్రకారం $\Gamma
      \Entails !B$, $\Delta \Entails !C$. మొత్తం
      \tetoken{వ్యుత్పత్తిలోని}{!!{derivation}}
      \tetoken{ఉపసంహరించని}{!!{undischarged}}
      పరికల్పనలు $\Gamma$, $\Delta$ కలిపినవే
      కాబట్టి $\Gamma \cup \Delta \Entails !B \land !C$
      అని చూపాలి.
      \sourcecorrection{OLTEINTSND-001}{మూలం ఈ సంయోగ పరిచయ కేసు లక్ష్యాన్ని $!A \land !B$గా వ్రాసింది; పూర్వాపేక్షలు, చివరి సత్య నిబంధన రెండూ $!B \land !C$ను నిర్ణయిస్తాయి కాబట్టి దానికి సరిచేశాం.}
      $\mSat{M}{\Gamma
        \cup \Delta}[w]$ అనుకుందాం. అప్పుడు
      $\mSat{M}{\Gamma}[w]$ కూడా. $\Gamma
      \Entails !B$ కాబట్టి $\mSat{M}{!B}[w]$.
      అదేవిధంగా $\mSat{M}{!C}[w]$. కనుక
      $\mSat{M}{!B \land !C}[w]$.}

  \item \tetoken{వ్యుత్పత్తి}{!!{derivation}} చివర
    $\Elim{\land}$ ఉంటే:
    \iftag{probAnd}{అభ్యాసం.}{$\Gamma$ అనే
      \tetoken{ఉపసంహరించని}{!!{undischarged}}
      పరికల్పనల నుంచి $!B
        \land !C$ పూర్వాపేక్ష యొక్క
      \tetoken{వ్యుత్పత్తి}{!!{derivation}},
      $\Gamma \Proves !B \land !C$ అని చూపుతుంది.
      ఆగమన పరికల్పన ప్రకారం $\Gamma \Entails !B
      \land !C$. $\Gamma \Entails !B$ అని
      చూపాలి. $\mSat{M}{\Gamma}[w]$ అనుకుందాం.
      $\Gamma
        \Entails !B \land !C$ కాబట్టి
      $\mSat{M}{!B \land !C}[w]$. అప్పుడు
      $\mSat{M}{!B}[w]$ కూడా. అదేవిధంగా
      $\Elim\land$ చివర~$!C$ వస్తే $\Gamma
      \Entails !C$.}

  \item \tetoken{వ్యుత్పత్తి}{!!{derivation}} చివర
    $\Intro{\lor}$ ఉంటే:
    \iftag{probOr}{అభ్యాసం.}{పూర్వాపేక్ష $!B$ అని,
      $!B$తో ముగిసే \tetoken{వ్యుత్పత్తిలో}{!!{derivation}}
      \tetoken{ఉపసంహరించని}{!!{undischarged}}
      పరికల్పనలు $\Gamma$ అని అనుకుందాం. అప్పుడు
      $\Gamma \Proves !B$; ఆగమన పరికల్పన ప్రకారం
      $\Gamma \Entails !B$. $\Gamma \Entails !B
      \lor !C$ అని చూపాలి. $\mSat{M}{\Gamma}[w]$
      అనుకుందాం. $\Gamma \Entails !B$ కాబట్టి
      $\mSat{M}{!B}[w]$. అప్పుడు $\mSat{M}{!B
        \lor !C}[w]$ కూడా. అదేవిధంగా పూర్వాపేక్ష~$!C$
      అయితే $\Gamma \Entails !C$; ఆ కేసులోనూ
      కోరిన వియోజనం వస్తుంది.}

  \item \tetoken{వ్యుత్పత్తి}{!!{derivation}}
    చివర~$\Elim{\lor}$ ఉంటే:
    \iftag{probOr}{అభ్యాసం.}{పూర్వాపేక్షల్లో మొదటిది
      \tetoken{ఉపసంహరించని}{!!{undischarged}}
      పరికల్పనలు~$\Gamma$ నుంచి $!B \lor !C$,
      రెండవది \tetoken{ఉపసంహరించని}{!!{undischarged}}
      పరికల్పనలు $\Delta_1 \cup \{!B\}$ నుంచి $!D$,
      మూడవది \tetoken{ఉపసంహరించని}{!!{undischarged}}
      పరికల్పనలు~$\Delta_2 \cup \{!C\}$ నుంచి
      $!D$కు చెందిన \tetoken{వ్యుత్పత్తులు}{!!{derivation}s}.
      కాబట్టి $\Gamma \Proves
      !B \lor !C$, $\Delta_1 \cup \{!B\} \Proves !D$,
      $\Delta_2
      \cup \{!C\} \Proves !D$. ఆగమన పరికల్పన
      ప్రకారం $\Gamma
      \Entails !B \lor !C$, $\Delta_1 \cup \{!B\}
      \Entails !D$, $\Delta_2 \cup \{!C\}
      \Entails !D$. $\Gamma \cup \Delta_1 \cup
      \Delta_2 \Entails !D$ అని నిరూపించాలి.

    $\mSat{M}{\Gamma \cup \Delta_1 \cup \Delta_2}[w]$
    అనుకుందాం. అప్పుడు $\mSat{M}{\Gamma}[w]$;
    $\Gamma \Entails !B \lor !C$ కాబట్టి
    $\mSat{M}{!B \lor !C}[w]$. $\mSat{M}{}$
    నిర్వచనం ప్రకారం $\mSat{M}{!B}[w]$ లేదా
    $\mSat{M}{!C}[w]$. అందువల్ల రెండు కేసులు:
    (a)~$\mSat{M}{!B}[w]$. అప్పుడు
    $\mSat{M}{\Delta_1 \cup \{!B\}}[w]$.
    \sourcecorrection{OLTEINTSND-002}{మూలం (a) కేసులో $[w]$ను గణిత-సంకేతం బయట ఉంచింది; నిర్వచిత లోక-సత్య సంకేతంలోకి చేర్చాం.}
    $\Delta_1 \cup \{!B\} \Entails !D$ కాబట్టి
    $\mSat{M}{!D}[w]$.
    \sourcecorrection{OLTEINTSND-003}{మూలంలోని రెండు అనుగమన వ్యక్తీకరణల్లో $!B$, $!C$లను సమితి-సమ్మేళనానికి నేరుగా కలిపింది; పూర్వ వ్యుత్పత్తి, ఆగమన పరికల్పనలకు అనుగుణంగా ఒక్కో మూలక సమితి జతచేశాం.}
    (b)~$\mSat{M}{!C}[w]$. అప్పుడు
    $\mSat{M}{\Delta_2 \cup \{!C\}}[w]$.
    $\Delta_2 \cup \{!C\} \Entails !D$ కాబట్టి
    $\mSat{M}{!D}[w]$. ఏ కేసులోనైనా
    కోరినట్లుగా $\mSat{M}{!D}[w]$.}

  \item \tetoken{వ్యుత్పత్తి}{!!{derivation}}
    చివర $\Intro{\lif}$ ద్వారా~$!B\lif !C$
    వస్తే, పూర్వాపేక్ష~$!C$. దాని
    \tetoken{వ్యుత్పత్తిలోని}{!!{derivation}}
    \tetoken{ఉపసంహరించని}{!!{undischarged}}
    పరికల్పనలు~$\Gamma \cup \{!B\}$. కాబట్టి
    $\Gamma \cup \{!B\} \Proves !C$; ఆగమన
    పరికల్పన ప్రకారం $\Gamma \cup \{!B\}
    \Entails !C$. $\Gamma \Entails !B \lif !C$
    అని చూపాలి.

    $\mSat{M}{\Gamma}[w]$ అనుకుందాం. $Rww'$
    అయ్యే ప్రతి $w'$ వద్ద $\mSat{M}{!B}[w']$
    అయితే $\mSat{M}{!C}[w']$ అని చూపాలి.
    కాబట్టి $Rww'$ మరియు $\mSat{M}{!B}[w']$
    అనుకుందాం. \olref[sem][rel]{prop:true-monotonic}
    ప్రకారం $\mSat{M}{\Gamma}[w']$.
    $\Gamma \cup \{!B\} \Entails !C$ కాబట్టి
    కోరినట్లుగా $\mSat{M}{!C}[w']$.

  \item \tetoken{వ్యుత్పత్తి}{!!{derivation}}
    చివర $\Elim{\lif}$ ద్వారా~$!C$ వస్తే,
    పూర్వాపేక్షలు $!B \lif !C$, $!B$.
    వాటి \tetoken{వ్యుత్పత్తులు}{!!{derivation}s}
    వరుసగా \tetoken{ఉపసంహరించని}{!!{undischarged}}
    పరికల్పనలు $\Gamma$, $\Delta$ నుంచి వస్తాయి.
    కాబట్టి $\Gamma \Proves !B \lif !C$,
    $\Delta
    \Proves !B$. ఆగమన పరికల్పన ప్రకారం
    $\Gamma \Entails !B \lif !C$, $\Delta
    \Entails !B$. $\Gamma \cup \Delta \Entails !C$
    అని చూపాలి.

    $\mSat{M}{\Gamma \cup \Delta}[w]$ అనుకుందాం.
    $\mSat{M}{\Gamma}[w]$ మరియు
    $\Gamma \Entails !B \lif !C$ కాబట్టి
    $\mSat{M}{!B \lif !C}[w]$. నిర్వచనం ప్రకారం
    $Rww'$ అయ్యే ప్రతి $w'$ వద్ద
    $\mSat{M}{!B}[w']$ అయితే
    $\mSat{M}{!C}[w']$. $R$ స్వావర్తనం కాబట్టి
    $Rww'$ అయ్యే $w'$లలో $w$ కూడా ఉంది; అంటే
    $\mSat{M}{!B}[w]$ అయితే
    $\mSat{M}{!C}[w]$. $\mSat{M}{\Delta}[w]$
    మరియు $\Delta
    \Entails !B$ కాబట్టి $\mSat{M}{!B}[w]$.
    కనుక కోరినట్లుగా $\mSat{M}{!C}[w]$.

  \item \tetoken{వ్యుత్పత్తి}{!!{derivation}}
    చివర $\FalseInt$ ద్వారా~$!A$ వస్తే,
    పూర్వాపేక్ష $\lfalse$. ఆ పూర్వాపేక్ష
    \tetoken{వ్యుత్పత్తిలోని}{!!{derivation}}
    \tetoken{ఉపసంహరించని}{!!{undischarged}}
    పరికల్పనలు~$\Gamma$. అప్పుడు $\Gamma
    \Proves \lfalse$. ఆగమన పరికల్పన ప్రకారం
    $\Gamma \Entails \lfalse$. $\Gamma
    \Entails !A$ అని చూపాలి.

    పరోక్షంగా నిరూపిద్దాం. $\Gamma \Entails/ !A$
    అయితే, ఏదో నమూనా~$\mModel{M}$, లోకం~$w$లకు
    $\mSat{M}{\Gamma}[w]$ మరియు
    $\mSat/{M}{!A}[w]$. కానీ $\Gamma \Entails
    \lfalse$ కాబట్టి $\mSat{M}{\lfalse}[w]$.
    నిర్వచనం ప్రకారం $\mSat/{M}{\lfalse}[w]$
    కావడంతో ఇది అసాధ్యం. కనుక $\Gamma
    \Entails !A$.
  \item \tetoken{వ్యుత్పత్తి}{!!{derivation}}
    చివర $\Intro\lnot$ ఉంటే: అభ్యాసం.
  \item \tetoken{వ్యుత్పత్తి}{!!{derivation}}
    చివర $\Elim\lnot$ ఉంటే: అభ్యాసం.
  \end{enumerate}
\end{proof}

\begin{prob}
  \olref[int][sc][snd]{thm:soundness} నిరూపణను
  పూర్తి చేయండి. $\Intro\lnot$, $\Elim\lnot$
  కేసులకు \olref[int][sem][rel]{defn:true-at-w}లోని
  $\mSat{M}{\lnot !A}[w]$ నిర్వచనాన్ని వాడండి;
  అంటే $\lnot !A$ను $!A \lif \lfalse$ ద్వారా
  నిర్వచించబడినదిగా పరిగణించవద్దు.
\end{prob}

\begin{prob}
  కింది \tetoken{సూత్రాలు}{!!{formula}s}
  అంతఃప్రజ్ఞావాద తర్కంలో
  \tetoken{వ్యుత్పాదించదగినవి}{!!{derivable}} కావని చూపండి:
  \begin{enumerate}
    \item $(!A \lif !B) \lor (!B \lif !A)$
    \item $(\lnot\lnot !A \lif !A) \lif (!A \lor \lnot !A)$
    \item $(!A \lif !B \lor !C) \lif \bigl((!A \lif !B)\lor(!A \lif !C)\bigr)$
  \end{enumerate}
\end{prob}

\end{document}
